\chapter{Berlioz, Hector} % (fold)
\label{cha:berlioz}

\section{Orchestral}

\noindent{Waverley (1828), Op. 1} \textit{LSO, Sir Colin Davis}\marginnote{high} \\
\noindent{Grande Ouverture des franc-juges (1826-1833), Op. 3} \textit{Les Si\`{e}cles, Fran\c{c}ois-Xavier Roth}\marginnote{high}  \\
\noindent{Le roi Lear (1831), Op. 4} \textit{LSO, Sir Colin Davis}\marginnote{high} \\
\noindent{Grande messe des morts (1837), Op. 5, H. 75} \textit{Robert Murray, Gabrieli Consort and Players, Wroclaw Philharmonic Choir \& Orchestra, Chetham's School of Music Brass Ensemble, Paull McCreesh}\marginnote{11/2011, max} \\  
\noindent{Les Nuits d'\'{e}t\'{e} (1840-1841), Op. 7} \textit{Marie-Nicole Lemieux, Orchestre Philharmonique de Monte Carlo, Kazuki Yamada}\marginnote{max}  \\
\noindent{Rêverie et caprice (Romance for Violin and Orchestra) (1841), Op. 8} \textit{James Ehnes, Melbourne SO, Andrew Davis}\marginnote{high} \\
\noindent{Le Carnaval romain (1844), Op. 9} \textit{LSO, Sir Colin Davis}\marginnote{high} \\
\noindent{Symphonie fantastique (1830), Op. 14} \textit{Anima Eterna, Jos van Immerseel}\marginnote{high}  \\
\noindent{Grande symphonie funèbre et triomphale (1840), Op. 15} \\
\noindent{Harold en Italie (1834), Op. 16} \textit{James Ehnes, Melbourne SO, Andrew Davis}\marginnote{high}  \\
\noindent{Roméo et Juliette (1839), Op. 17} \textit{LSO, Sir Colin Davis}\marginnote{high}  \\
\noindent{Tristia (1831-1844), Op. 18, H. 119} \textit{LSO, Sir Colin Davis}\marginnote{high} \\  
\noindent{Le Corsaire (1844), Op. 21} \textit{LSO, Sir Colin Davis}\marginnote{high} \\
\noindent{Benvenuto Cellini (1836-1838), Op. 23} \textit{LSO, Sir Colin Davis}\marginnote{high} \\ 
\noindent{Rob Roy (1831)} \textit{Melbourne SO, Andrew Davis}\marginnote{high} \\
\noindent{B\'{e}atrice et B\'{e}n\'{e}dict (1860-1862), H. 138} \textit{LSO, Sir Colin Davis}\marginnote{high} \\ 
\noindent{Marche Troyenne (1864)} \\



